% =====================================================================
% 参考文献
%
% 用 BibTeX + natbib 生成，样式为 gbt7714-numerical.bst（GB/T 7714-2015，
% 连续编号、方括号、上标引用）。条目写在根目录的 refs.bib 里。
%
% 编译顺序（BibTeX 不能省）：
%     xelatex main  →  bibtex main  →  xelatex main  →  xelatex main
% 在 VS Code / LaTeX Workshop 下选 “xe->bib->xe->xe” 配方即可。
%
% 正文里引用：\upcite{键名}（上标方括号，符合 GB/T 7714 与附件 3 示例），
% 需要平排方括号时用 \cite{键名}。
%
% 本页排版说明
%   * 标题“参考文献”由本页的 \bibsection（重定义为文档类的章标题格式）排出，
%     同时手工写一条目录项——因为这里不走 \chapter，latex 不会自动加目录。
%   * 条号顶格、续行缩进 2 字符（文档类已显式设置数字制列表缩进），
%     与附件 3 的版式一致。
% =====================================================================

\clearpage
\phantomsection
\addcontentsline{toc}{chapter}{参考文献}%
\begingroup
\setstretch{1.5}%
\songti\zihao{-4}%
% 对齐附件 3：取消文献条目间额外间距，正文基线约 23.3pt。
\setlength{\bibsep}{0pt}%
\setlength{\baselineskip}{23.3pt}%
\makeatletter
\let\@bibsetup\NAT@bibsetnum
\renewcommand{\bibsection}{%
  \vspace*{-3.386pt}%
  \ujs@chapterfmt{\bibname}%
  \vspace{-5.815pt}%
}
\makeatother
\bibliography{refs}
\endgroup





